\section{Parameter des Prototyps}
\label{anh:parameter}

Die Tabelle nennt die im Quellcode des Prototyps hinterlegten Vorgaben, auf die
\autoref{sec:Entwurf} verweist, Stand 15. Oktober 2026, konfigurierbare Werte mit ihrem Standardwert
und, wo die Integrationsumgebung abweicht, mit deren Wert. Metriken gehen per \ac{OTLP}-Push an die Plattform.

\begin{table}[H]
  \caption{Parameter des Prototyps nach Komponente}
  \label{tab:parameter}
  \scriptsize
  \renewcommand{\arraystretch}{1.0}
  \begin{tabularx}{\textwidth}{@{}>{\RaggedRight\arraybackslash}p{0.13\textwidth}
    >{\RaggedRight\arraybackslash}p{0.30\textwidth} >{\RaggedRight\arraybackslash}X@{}}
    \hline
    Komponente & Parameter & Wert \\
    \hline
    Runtime & \texttt{deno\_core} & Version 0.412, exakt festgelegt \\
    & Node-Quellcode der Vendor-Bundles & Version 26.10.0, per Commit festgelegt \\
    & Laufzeit-Snapshot & rund 1,6\,MB \\
    & lazy geladene Bundles & Node-Module rund 690\,KB, Web Streams rund 480\,KB \\[2pt]
    Worker & Slots je Pod & 4 \\
    & Heap je Isolate & Pod-Limit / (2 $\times$ Slots), mindestens 64\,MB, ebenso viel als Limit für Puffer außerhalb des Heaps \\
    & Headroom nach Erreichen des Heap-Limits & 32\,MB \\
    & Frist synchron / Hintergrund & 50\,s / 830\,s \\
    & Body inline / maximal & 64\,KiB / 100\,MiB \\
    & Bundle entpackt, maximal & 64\,MB \\
    & Bundle-Cache & 128\,MB je Prozess, Schlüssel SHA-256, Verdrängung des am
      längsten nicht genutzten Eintrags \\
    & Selbstbeendigung ohne Auftrag (Rückfall) & 5\,min (Umgebungsvariable), geprüft alle 30\,s, nur im
      zugewiesenen Zustand und nur zwischen zwei Pulls \\
    & Mandanten-Consumer & dauerhaft, Repliken wie der Stream (drei), höchstens vier Zustellungen, BackOff 5, 10 und
      20\,s, gelöscht nach 300\,s ohne Nutzung, offene Zustellungen höchstens Slots $\times$ 10,
      ein Pull je Nachricht mit Frist 2\,s, Warten auf den Leader bis 2\,s \\
    & Lesen des Consumers & Frist 2\,s, drei Versuche, Reconnect nach einem Versuch ohne Antwort \\
    & NATS-Verbindung bei der ersten Zuweisung & 500\,ms je Versuch, Wiederholung bis 30\,s, danach höchstens 66
      Reconnects, Wartezeit bis 5\,s \\
    & Bundle-Abruf & 3 Versuche (50\,ms, 200\,ms) bei Transportfehlern, 5xx und 429 \\
    & Callback und Log-Upload & 3 Versuche (500\,ms, 1\,s) bei Transportfehlern und 5xx, je 30\,s Frist \\
    & Metriken & \ac{OTLP}-Push alle 10\,s, Buckets 1\,ms bis 30\,s \\
    & Readiness-Probe & \texttt{/ready}, alle 10\,s \\[2pt]
    Gateway & Wartezeit synchron & 53\,s, unter den 55\,s des Edge-Proxys \\[2pt]
        Control Plane & Pool-Zielgröße & 3 (Code), 40 (Integrationsumgebungen) \\
    & Kommando an einen Worker (assign, unassign) & \ac{HTTP}-\texttt{POST} an den Transport des
      Pods, Antwort in der Response, Frist 10\,s \\
    & Reserve für Erweiterungen & Zielgröße / 2 \\
    & Rotation & nach 1\,h Zuweisungsdauer, geprüft alle 10\,s \\
    & Rückzug ohne Auftrag & \emph{retire} nach 1\,min ohne Nachricht des Mandanten, geprüft alle 10\,s,
      überzählige Pods über dem Maximum der letzten Minute \\
    & Pool-Vorhersage & Nachfrage je Minute, Rückblick vier Wochen und sieben Tage, Fenster 3\,min mit
      Vorlauf, Sicherheitsfaktor 1,3, Haltedauer 3\,min, Obergrenze 50, für die Messungen aus \\
    & Pods ohne Eintrag & gelöscht nach 5\,min \\
    & Neuzuweisung nach Pod-Tod & sofort, falls Aufträge warten, Wiederholung nach 1, 2, 4, 8,
      16\,s \\
    & Anlegen beim Erstkontakt & höchstens eines je Mandant in 30\,s, Warten auf Bereitschaft
      bis 20\,s, geprüft je 1\,s \\
    & Watcher-Consumer & bis zu 20 Zustellungen \\
    & Lastskalierung je Mandant & Ziel-Auslastung 0,5 (Code) bzw.\ 0,75 (Integrationsumgebungen)
      der Slots, höchstens 10 Pods, mindestens
      1, Lastprüfung alle 250\,ms, erste Erweiterung frühestens 2\,s nach dem Erstkontakt, Abstand
      zwischen Erweiterungen 2\,s, Pool danach auf Zielgröße aufgefüllt \\
    & Metriken & \ac{OTLP}-Push alle 15\,s \\[2pt]
    Pod & Ressourcen Sandbox & \ac{CPU}-Request 10\,m ohne Limit, Speicher 32\,Mi Request und 512\,Mi
      Limit in der Integrationsumgebung, 128\,Mi Request im Code \\
    & Ressourcen Transport & Speicher 32\,Mi Request und 256\,Mi Limit \\
    & Token der Zuweisung & RS256, gültig 80\,min, also Zuweisungsdauer plus Grace Period und Reserve \\
    & Grace Period & 900\,s \\
    & Service Mesh & Sidecar-Injection deaktiviert \\
    & Egress & Worker-Pod: Internet nur über den Egress-Proxy auf Port 3128, NATS über ein Label der Cluster-Richtlinien, kein
      Redis. Proxy-Pod: jeder öffentliche Port, private Bereiche ausgenommen \\
    \hline
  \end{tabularx}
\end{table}

\section{Header-Filter und Pod-Zustände}
\label{anh:header-zustaende}

\autoref{tab:header} nennt die Header, die der Worker nach \autoref{sec:bausteine} vor dem Skript
verbirgt oder aus dessen Antwort entfernt. \autoref{tab:podzustaende} führt die Zustände eines
Worker-Pods aus \autoref{sec:zuteilung} mit den Ereignissen, die sie beenden.

\begin{table}[H]
  \caption{Vom Worker gefilterte Header}
  \label{tab:header}
  \small
  \begin{tabularx}{\textwidth}{@{}>{\RaggedRight\arraybackslash}p{0.26\textwidth}
    >{\RaggedRight\arraybackslash}X@{}}
    \hline
    Richtung & Header \\
    \hline
    Request, vor dem Skript verborgen & Signaturen der Plattform, Identität des Service Mesh,
      Proxy-Header der Zwischenschritte \\[2pt]
    Response, verworfen & Cookies sowie \ac{CORS}, \ac{CSP}, \ac{HSTS} und der Schutz gegen
      Framing, mit denen ein Skript die Richtlinien der Plattform-Origin schwächen könnte \\[2pt]
    \texttt{fetch}, bei Redirects & Header, die den Aufrufer ausweisen, folgen keinem Redirect zu
      einer anderen Origin \\
    \hline
  \end{tabularx}
\end{table}

\begin{table}[H]
  \caption{Zustände eines Worker-Pods und ihre Übergänge}
  \label{tab:podzustaende}
  \small
  \begin{tabularx}{\textwidth}{@{}l >{\RaggedRight\arraybackslash}X
    >{\RaggedRight\arraybackslash}X@{}}
    \hline
    Zustand & Bedeutung & Übergänge \\
    \hline
    angelegt & über die Kubernetes-\acs{API} erzeugt, noch nicht bereit, in keiner der drei
      Mengen & bei Bereitschaft \emph{frei}, ohne Bereitschaft nach der Schonfrist aus
      \autoref{tab:parameter} gelöscht \\[2pt]
    frei & bereit, im Pool & beim Erstkontakt in einem atomaren Schritt nach
      \emph{reserviert} \\[2pt]
    reserviert & Zuweisung läuft, Kommando über \ac{HTTP}, Antwort innerhalb der Frist & bei
      Bestätigung \emph{zugewiesen}, bei einem fehlgeschlagenen Schritt zurück nach \emph{frei} \\[2pt]
    zugewiesen & genau einem Mandanten zugeordnet, nie wieder im Pool & Rotation nach Ablauf
      der Zuweisungsdauer, Rückzug ohne Auftrag oder Entzug des Knotens, jeweils nach
      \emph{beendet}, nie zurück nach \emph{frei} \\[2pt]
    beendet & Pod hat einen \texttt{DeletionTimestamp} bekommen oder ist gestorben, ein Informer
      meldet es & aus jedem Zustand. War es der letzte Pod
      eines Mandanten und warten Aufträge, Neuzuweisung sofort, Wiederholung mit wachsendem
      Abstand nach \autoref{tab:parameter} \\
    \hline
  \end{tabularx}
\end{table}

Nach einem Neustart liest die Control Plane den Zustand aus Redis zurück, entfernt zuerst Pods, die
es nicht mehr gibt, und wiederholt offene Zuweisungen. Das ist möglich, weil der Worker eine
wiederholte Zuweisung an denselben Mandanten als Bestätigung behandelt.


\section{Codeauszüge}
\label{sec:codeauszuege}

Die Auszüge stammen aus dem Prototyp. Sie sind gekürzt, die Kommentare sind übersetzt oder entfallen.

\lstdefinelanguage{GoLang}{
  keywords={func, return, if, nil, var, const, type, struct, string, error, package, import, range},
  sensitive=true, comment=[l]{//}, morestring=[b]"
}
\lstdefinelanguage{RustLang}{
  keywords={fn, pub, impl, match, let, if, else, self, Self, true, false, enum, struct, use, mut, const},
  sensitive=true, comment=[l]{//}, morestring=[b]"
}

\begin{lstlisting}[language=GoLang, caption={Rechte der NATS-Verbindung eines Workers (Auth Callout, \texttt{WorkerGrant.Permissions}, Go)}, label={lst:callout-rechte}]
func (g WorkerGrant) Permissions() Permissions {
	consumer := Stream + ".scriptworker-tenant-" + g.TenantID
	return Permissions{
		PubAllow: []string{
			"$JS.API.CONSUMER.INFO." + consumer,
			"$JS.API.CONSUMER.MSG.NEXT." + consumer,
			"$JS.ACK." + consumer + ".>",
			// Antworten nur an Aufrufer des eigenen Mandanten
			"dv.processstudio.replies.*." + g.TenantID + ".>",
		},
		PubDeny: []string{
			"$JS.API.CONSUMER.CREATE.>",
			"$JS.API.CONSUMER.DURABLE.CREATE.>",
			"$JS.API.CONSUMER.DELETE.>",
			"$JS.API.STREAM.>",
			"$JS.API.DIRECT.GET.>",
			"dv.processstudio.cmds.>",
		},
		SubAllow: []string{"_INBOX_" + g.WorkerID + ".>"},
		SubDeny:  []string{"_INBOX.>", "dv.processstudio.>", "$JS.>"},
	}
}
\end{lstlisting}

\begin{lstlisting}[language=GoLang, caption={Prüfung der Verbindung eines Workers im Auth Callout (\texttt{Service.Authorize}, Go)}, label={lst:callout-token}]
claims, err := s.tokens.Verify(ctx, c.Token)
if err != nil {
	return domain.WorkerGrant{}, domain.Rejection{Reason: "invalid assignment token: " + err.Error()}
}
if !domain.ValidWorkerID(claims.Subject) || claims.Subject != c.User {
	return domain.WorkerGrant{}, domain.Rejection{Reason: "the token is for another worker"}
}
if !domain.ValidTenantID(claims.TenantID) {
	return domain.WorkerGrant{}, domain.Rejection{Reason: "the token's tenant is not a valid tenant id"}
}
if !identityMatches(c.Password, claims.IdentityHash) {
	return domain.WorkerGrant{}, domain.Rejection{Reason: "the identity secret doesn't match the token"}
}
return domain.WorkerGrant{WorkerID: claims.Subject, TenantID: claims.TenantID, ExpiresAt: claims.ExpiresAt}, nil
\end{lstlisting}

\begin{lstlisting}[language=RustLang, caption={Absender im eigenen Pod werden als Quelle von Kommandos abgelehnt (\texttt{CommandHandler::is\_own\_pod}, Rust)}, label={lst:eigener-pod}]
fn is_own_pod(&self, peer: IpAddr) -> bool {
    let peer = match peer {
        IpAddr::V6(v6) => v6.to_ipv4_mapped().map_or(peer, IpAddr::V4),
        v4 => v4,
    };
    peer.is_loopback() || Some(peer) == self.pod_ip
}
\end{lstlisting}

\begin{lstlisting}[language=RustLang, caption={Zuordnung des Profils zu den Bindings (\texttt{BindingSet::for\_profile}, Rust)}, label={lst:bindings}]
pub fn for_profile(profile: Profile) -> Self {
    match profile {
        Profile::Webhook => BindingSet {
            fetch: true,
            net: true,
            tls: true,
            console_capture: true,
            node_globals: true,
        },
        // Agentenprofil: erzeugter, möglicherweise beeinflusster Quelltext,
        // bekommt nie Socket-Bindings, was er auch importiert.
        Profile::Agent => BindingSet {
            fetch: true,
            net: false,
            tls: false,
            console_capture: false,
            node_globals: false,
        },
    }
}
\end{lstlisting}


\section{Weitere Messdaten zu Kapitel 5}
\label{sec:messdaten}

\begin{table}[H]
  \caption{Stand des Prototyps je Beleg}
  \label{tab:stand}
  \small
  \begin{tabularx}{\textwidth}{@{}>{\RaggedRight\arraybackslash}X >{\RaggedRight\arraybackslash}p{0.12\textwidth} >{\RaggedRight\arraybackslash}p{0.27\textwidth}@{}}
    \hline
    Beleg & Datum & Stand \\
    \hline
    S1a, S1b, S2, S3, S4, M5, Zustellung 06 bis 08 und Absturz, T1 bis T9, SandboxEval, M1 und M3 & 29.\ und 30.\ Oktober & Endstand: Gateway \texttt{b2aa8e7}, Executor \texttt{5db344b}, Control Plane \texttt{55331df} \\
    Versuch mit einem Slot je Pod & 30.\ Oktober & wie Endstand, ein Slot je Pod \\
    Noisy Neighbour, Lauf P1 und K1 & 2.\ November, morgens & wie Endstand \\
    Noisy Neighbour, Läufe P2 bis K3, Proxy-Abmeldung im Hintergrund, Mikrobenchmarks, Prozessstart unter gVisor, reale Bundles & 2.\ November & Gateway \texttt{e52c21d}, Executor \texttt{5276a03}, Control Plane 256\,MiB \\
    S1a und S2 vor dem Endstand, M5 mit Prober, T7 und T9 & 26.\ bis 28.\ Oktober & ohne Worker-ID in der Antwort, ab 26.\ Oktober mit Transport und Sandbox, ab 28.\ Oktober Proxy in enforce \\
    Kaltstartverlauf, S1b und S2 Läufe 28 und 30, Pod-Ressourcen & bis 15.\ Oktober & vor der Trennung in Transport und Sandbox \\
    T10 & 28.\ Oktober, lokal & mit Auth Callout \\
    \hline
  \end{tabularx}
\end{table}

\begin{table}[H]
  \caption{Kaltstart je Mandant, Einzelläufe vom 29.\ Oktober, Reihenfolge P1, K1, P2, K2, P3, K3, dann Pool 0, Fehler (502 \enquote{no route to host}) zählen mit ihrer Dauer und nicht als unter zwei Sekunden}
  \label{tab:s1-laeufe}
  \footnotesize
  \setlength{\tabcolsep}{4pt}
  \begin{tabularx}{\textwidth}{@{}l r r r r r r@{}}
    \hline
    Lauf & Median & P95 & P99 & Maximum & unter 2\,s & Fehler \\
    \hline
    \multicolumn{7}{@{}l}{\emph{verteilt, 90 Mandanten, einer je Sekunde}} \\
    P1 Pool 40 & 281\,ms & 411\,ms & 529\,ms & 555\,ms & 90 & 0 \\
    P2 Pool 40 & 278\,ms & 527\,ms & 648\,ms & 1,23\,s & 90 & 0 \\
    P3 Pool 40 & 308\,ms & 519\,ms & 1,27\,s & 1,37\,s & 90 & 0 \\
    K1 Ausgang & 2,54\,s & 3,18\,s & 3,26\,s & 3,31\,s & 0 & 0 \\
    K2 Ausgang & 2,76\,s & 3,30\,s & 4,69\,s & 5,03\,s & 0 & 0 \\
    K3 Ausgang & 2,64\,s & 3,08\,s & 4,82\,s & 5,19\,s & 3 & 0 \\
    P0-1 Pool 0 & 2,52\,s & 6,62\,s & 7,39\,s & 7,59\,s & 36 & 0 \\
    P0-2 Pool 0 & 2,54\,s & 7,76\,s & 12,8\,s & 23,5\,s & 34 & 0 \\
    P0-3 Pool 0 & 2,65\,s & 8,15\,s & 9,62\,s & 9,65\,s & 29 & 0 \\
    \multicolumn{7}{@{}l}{\emph{Schub, 90 Mandanten zugleich}} \\
    P1 Pool 40 & 3,53\,s & 7,68\,s & 8,79\,s & 10,8\,s & 39 & 0 \\
    P2 Pool 40 & 4,25\,s & 8,57\,s & 10,9\,s & 11,0\,s & 40 & 0 \\
    P3 Pool 40 & 3,42\,s & 7,36\,s & 7,91\,s & 8,72\,s & 40 & 0 \\
    K1 Ausgang & 8,23\,s & 15,3\,s & 20,2\,s & 20,2\,s & 0 & 0 \\
    K2 Ausgang & 8,12\,s & 13,9\,s & 21,8\,s & 21,8\,s & 0 & 5 \\
    K3 Ausgang & 8,64\,s & 16,5\,s & 23,6\,s & 24,3\,s & 0 & 3 \\
    P0-1 Pool 0 & 6,55\,s & 11,3\,s & 12,5\,s & 12,6\,s & 0 & 0 \\
    P0-2 Pool 0 & 7,23\,s & 11,5\,s & 12,6\,s & 14,5\,s & 0 & 0 \\
    P0-3 Pool 0 & 7,65\,s & 11,9\,s & 13,1\,s & 13,9\,s & 0 & 0 \\
    \hline
  \end{tabularx}
\end{table}

\begin{table}[H]
  \caption{Verlauf des Overheads im Lastprofil bis zum 28.\ Oktober in Millisekunden, k6, je Lauf rund 33 Minuten, Fehler mit Status 502 nicht in den Perzentilen}
  \label{tab:overhead-verlauf}
  \footnotesize
  \begin{tabularx}{\textwidth}{@{}l >{\RaggedRight\arraybackslash}X r r r r r r r@{}}
    \hline
    Lauf & Modell, Modus & Aufrufe & Fehler & P50 & P90 & P95 & P99 & Max \\
    \hline
    18 & Ausgangszustand, Profil & 56.976 & 408 & 259 & 1.190 & 2.310 & 4.570 & 50.530 \\
    20 & Prototyp, Wide & 68.881 & 1 (client) & 131 & 257 & 306 & 373 & 3.960 \\
    28 & Prototyp, Profil & 57.001 & 0 & 127 & 164 & 180 & 268 & 15.460 \\
        30 & Prototyp, Profil, mit Reconnect & 57.001 & 8 (7 client) & 127 & 163 & 178 & 276 & 5.580 \\
    26.10. & Prototyp nach der Trennung in Transport und Sandbox, Profil & 57.001 & 0 & 128 & 165 & 183 & 434 & 3.404 \\
        26.10. & Prototyp, Control Plane beendet den Leerlauf, Profil & 57.001 & 14 (502) & 138 & 239 & 293 & 405 & 3.443 \\
    28.10. & Ausgangszustand, Profil, Ingress-Limit angehoben & 56.989 & 6 (502) & 301 & 490 & 746 & 3.410 & 16.500 \\
    28.10. & Prototyp mit Egress-Proxy, Profil & 57.001 & 0 & 155 & 253 & 296 & 399 & 12.260 \\

    \hline
  \end{tabularx}
\end{table}

\begin{table}[H]
  \caption{Ressourcen unter dem Lastprofil, Spitzen über 33 Minuten, 29.\ Oktober, cAdvisor meldet unter gVisor nur Pods}
  \label{tab:ressourcen-endstand}
  \footnotesize
  \setlength{\tabcolsep}{4pt}
  \begin{tabularx}{\textwidth}{@{}>{\RaggedRight\arraybackslash}X r r r r r r@{}}
    \hline
     & P1 & P2 & P3 & K1 & K2 & K3 \\
    \hline
    Pods der Laufzeit gleichzeitig & 207 & 202 & 202 & 391 & 382 & 376 \\
    Speicher der Laufzeit-Pods, GiB & 15,0 & 14,4 & 14,6 & 34,1 & 33,8 & 33,0 \\
    \acs{CPU} der Laufzeit-Pods, Mittel / Spitze, Kerne & 3,44 / 4,06 & 3,35 / 4,03 & 3,33 / 3,94 & 3,28 / 4,10 & 3,08 / 3,94 & 3,22 / 3,82 \\
    Spitzenspeicher je Pod, Median / P90 / Max, MiB & 70,5 / 81,7 / 103 & & & 82,7 / 103 / 133 & & \\
    Steuerung und Routing, Spitze, MiB & 267 & & & 2.299 & & \\
    Pool im Ruhezustand, 40 Worker zu 78,7\,MiB, MiB / Kerne & 3.111 / 1,02 & & & 0 & & \\
    \hline
  \end{tabularx}
\end{table}

\begin{table}[H]
  \caption{Noisy Neighbour, Latenz der Mandanten B beim Client, P50 / P95 in Millisekunden, 2.\ November, K3 mit 150 Kaltstarts vor der Last}
  \label{tab:b6}
  \footnotesize
  \setlength{\tabcolsep}{4pt}
  \begin{tabularx}{\textwidth}{@{}>{\RaggedRight\arraybackslash}X r r r r r r@{}}
    \hline
     & P1 & P2 & P3 & K1 & K2 & K3 \\
    \hline
    nur B, vorher & 235 / 324 & 236 / 390 & 217 / 268 & 314 / 494 & 319 / 536 & 352 / 1.723 \\
    A und B, Knoten ab 90\,\% & 421 / 898 & 327 / 422 & 378 / 469 & 506 / 660 & 488 / 666 & 475 / 620 \\
    A und B, Knoten 70 bis 90\,\% & 288 / 845 & 293 / 376 & 273 / 356 & 394 / 532 & -- & -- \\
    A und B, Knoten unter 70\,\% & 242 / 440 & 222 / 366 & 228 / 305 & 341 / 461 & 333 / 442 & 312 / 407 \\
    nur B, nachher & 227 / 306 & 226 / 369 & 212 / 266 & 324 / 447 & 320 / 485 & 319 / 474 \\
    Knoten ab 90\,\% & 1 & 2 & 2 & 9 & 7 & 9 \\
    Pods von A & 10 auf 7 Knoten & 10 auf 7 & 10 auf 7 & etwa 50 auf 12 & etwa 50 auf 12 & etwa 50 auf 12 \\
    Fehler & 0 & 2 (Client) & 0 & 7 & 5 & 50 (49 $\times$ 500) \\
    \hline
  \end{tabularx}
\end{table}

\begin{table}[H]
  \caption{Reale Bundles aus den häufigsten Paketen des Bestands, P50 / P95 in Millisekunden, 2.\ November, je Paket fünf Mandanten mit einem Kaltaufruf und 20 warmen Aufrufen}
  \label{tab:bundles}
  \footnotesize
  \begin{tabularx}{\textwidth}{@{}>{\RaggedRight\arraybackslash}X r r r r@{}}
    \hline
    Paket & Prototyp kalt & Prototyp warm & Ausgang kalt & Ausgang warm \\
    \hline
    \enquote{Hello World} & 405 / 627 & 168 / 490 & 2.727 / 2.771 & 302 / 654 \\
    \texttt{@dvelop-sdk/dms} & 464 / 659 & 175 / 253 & 2.758 / 2.833 & 326 / 822 \\
    \texttt{axios} & 547 / 644 & 243 / 338 & 3.234 / 5.280 & 390 / 665 \\
    \texttt{pdf-lib} & 579 / 813 & 213 / 327 & 2.766 / 3.161 & 430 / 752 \\
    \texttt{pdfmake} & 1.077 / 1.409 & 289 / 576 & 3.616 / 5.024 & 649 / 1.286 \\
    \hline
  \end{tabularx}
\end{table}

\begin{table}[H]
  \caption{Agentenprofil auf dem Prototyp, Endstand vom 2.\ November, k6 vom Laptop über den Edge-Proxy. Warm: fünf Mandanten mit zusammen 10 Aufrufen je Sekunde
    über drei Minuten, je 1.801 Aufrufe, Agent und Webhook-Profil (Hello World) direkt nacheinander, P50 / P95 / P99 in Millisekunden.
    Kaltstart: 30 Mandanten seit dem 30.\ Oktober ohne Aufruf, ein Erstaufruf je Sekunde, 30 Sekunden später je zehn warme Aufrufe}
  \label{tab:agent-latenz}
  \footnotesize
  \begin{tabularx}{\textwidth}{@{}>{\RaggedRight\arraybackslash}X r r@{}}
    \hline
    Messgröße & Agent & Webhook-Profil \\
    \hline
    Client & 85,0 / 121,9 / 140,8 & 118,0 / 183,0 / 238,6 \\
    Gateway gesamt & 23,9 / 32,8 / 49,9 & 49,3 / 107,7 / 215,3 \\
    Skript auflösen (Skriptdienst) & -- & 25,7 / 65,9 / 126,6 \\
    Registrierung beim Proxy & 7,8 / 15,7 / 16,8 & -- \\
    \texttt{execute script}, NATS bis Antwort & 11,9 / 16,4 / 26,2 & 12,5 / 26,6 / 33,1 \\
    Isolate erzeugen, P50 / P95 & 8 / 10 & -- \\
    Skript ausführen, P50 / P95 & 1 / 3 & -- \\
    \multicolumn{3}{@{}l}{\emph{Kaltstart}} \\
    Erstaufruf, Client, P50 / P95 / Maximum & 546 / 647 / 954 & -- \\
    Erstaufruf, Gateway, P50 / P95 / P99 & 157 / 852 / 1.074 & -- \\
    Warme Aufrufe danach, Client, P50 / P95 / Maximum & 84 / 121 / 174 & -- \\
    \hline
  \end{tabularx}
\end{table}

\begin{table}[H]
  \caption{Agentenprofil, Größe von Ein- und Ausgabe, ein Mandant, je Größe 20 Aufrufe nacheinander, Client P50 / P95 in Millisekunden. Oberhalb
    von 64\,KiB geht der Körper über S3 mit vorsignierten Adressen}
  \label{tab:agent-groesse}
  \footnotesize
  \setlength{\tabcolsep}{3pt}
  \begin{tabularx}{\textwidth}{@{}l r r r r r r r@{}}
    \hline
     & 1\,KiB & 32\,KiB & 63\,KiB & 65\,KiB & 256\,KiB & 1\,MiB & 4\,MiB \\
    \hline
    Eingabe & 105 / 124 & 100 / 156 & 134 / 177 & 237 / 392 & 371 / 549 & 510 / 1.127 & 1.102 / 1.937 \\
    Ausgabe & 85 / 118 & 91 / 131 & 102 / 138 & 155 / 337 & 231 / 393 & 335 / 611 & 843 / 1.392 \\
    \hline
  \end{tabularx}
\end{table}

\begin{table}[H]
  \caption{Agentenprofil, ausgehende Anfragen an das eigene System über den Proxy in enforce, je Anzahl 30 Läufe nacheinander, Millisekunden}
  \label{tab:agent-egress}
  \footnotesize
  \begin{tabularx}{\textwidth}{@{}r r r r@{}}
    \hline
    Anfragen je Lauf & Client P50 / P95 & Skript P50 & je Anfrage laut Proxy P50 / P95 \\
    \hline
    0 & 81 / 126 & 1 & -- \\
    1 & 176 / 240 & 95 & 75 / 108 \\
    5 & 544 / 599 & 454 & 74 / 110 \\
    10 & 981 / 1.098 & 882 & 74 / 109 \\
    \hline
  \end{tabularx}
\end{table}

\begin{table}[H]
  \caption{Agentenprofil, Parallelität auf einem Mandanten, geschlossene Schleife, je Stufe 60 Sekunden, Millisekunden. Alle Aufrufe gelangen mit
    Status 200 durch}
  \label{tab:agent-parallel}
  \footnotesize
  \begin{tabularx}{\textwidth}{@{}r r r r r@{}}
    \hline
    gleichzeitig & Aufrufe & Aufrufe/s & Client P50 / P95 / P99 & Gateway P50 / P95 / P99 \\
    \hline
    1 & 619 & 10,4 & 89 / 125 / 139 & 23,9 / 33,0 / 54,2 \\
    4 & 2.481 & 41,4 & 89 / 130 / 149 & 23,9 / 33,1 / 53,8 \\
    8 & 4.986 & 83,2 & 89 / 127 / 149 & 23,9 / 33,0 / 51,9 \\
    16 & 8.543 & 142,4 & 109 / 161 / 249 & 24,2 / 37,8 / 60,7 \\
    32 & 16.932 & 282,3 & 110 / 157 / 185 & 26,3 / 59,3 / 70,3 \\
    64 & 30.805 & 512,9 & 120 / 185 / 230 & 32,9 / 108,9 / 142,0 \\
    \hline
  \end{tabularx}
\end{table}

\begin{table}[H]
  \caption{Start eines Isolates und eines Node.js-Prozesses, P50 / P95 in Millisekunden}
  \label{tab:prozessstart}
  \footnotesize
  \begin{tabularx}{\textwidth}{@{}>{\RaggedRight\arraybackslash}X r@{}}
    \hline
    Messung & P50 / P95 \\
    \hline
    \multicolumn{2}{@{}l}{\emph{Cluster, Prozess je \texttt{spawnSync}, 40 Läufe, Node v24.22.1}} \\
    \texttt{node -e 0}, gVisor & 53,1 / 58,7 \\
    \texttt{node -e 0}, runc & 26,0 / 32,0 \\
    \texttt{-{}-snapshot-blob}, leerer Snapshot, gVisor / runc & 65,6 / 72,9 und 30,2 / 35,6 \\
    \texttt{-{}-snapshot-blob} mit \texttt{util}, \texttt{events}, \texttt{url}, gVisor / runc & 67,1 / 77,6 und 27,8 / 33,1 \\
    Isolate des Webhook-Profils bauen, warm, im Cluster & 18,7 / 77 \\
    \multicolumn{2}{@{}l}{\emph{Lokal, Mac (arm64), Release-Build des Executors \texttt{5db344b}, 300 Läufe}} \\
    Isolate bauen, Webhook-Profil / Agentenprofil & 3,30 / 3,61 und 2,64 / 2,96 \\
    \enquote{Hello World} im frischen Isolate, gesamt & 3,92 / 4,39 \\
    Prozess \texttt{/usr/bin/true} & 3,68 / 8,94 \\
    Prozess \texttt{node -e 0}, 50 Läufe & 68,9 / 77,6 \\
    \hline
  \end{tabularx}
\end{table}

\begin{table}[H]
  \caption{Herleitung von S6 aus der Bestandsanalyse und den Prüfungen der Kompatibilitätsschicht}
  \label{tab:s6}
  \small
  \begin{tabularx}{\textwidth}{@{}>{\RaggedRight\arraybackslash}X r >{\RaggedRight\arraybackslash}p{0.34\textwidth}@{}}
    \hline
    Gruppe & Mandanten & Wertung \\
    \hline
    Alle Mandanten & 1.718 & \\
    Ohne blockierendes Merkmal & 1.005 & läuft \\
    Nur ein Netzwerk-Binding blockiert (HTTP 629, Sockets 63, beide 44) & 648 & läuft über die Schicht, Annahme \\
    Übrige, nur \texttt{http}, \texttt{https}, \texttt{net} oder \texttt{tls} & 16 & läuft über die Schicht, Annahme \\
    Übrige, Skripte gelesen und im Probelauf geprüft & 49 & \\
    \quad schreiben nach \texttt{/tmp} & 34 & nicht schlechter, scheitert schon im Ausgangszustand \\
    \quad lesen Windows-Pfade & 4 & nicht schlechter, ebenso \\
    \quad Import ungenutzt & 4 & läuft \\
    \quad schreiben nach \texttt{\_\_dirname} & 2 & läuft in beiden \\
    \quad \texttt{child\_process} & 1 & absichtlich blockiert, der eine nicht lauffähige Mandant \\
    \quad gebündelte Skripte & 4 & nicht unterstützt, nicht als Ausfall gezählt \\
    S6, bis zu & 1.717 von 1.718 & 99,9\,\%, mit den vier gebündelten als Ausfall 1.713, 99,7\,\% \\
    \hline
  \end{tabularx}
\end{table}


\begin{table}[H]
  \caption{Unsicherheit der Kaltstartläufe aus \autoref{tab:kaltstart-vergleich}: konservatives 95-\%-Intervall je Lauf für P95 aus
    den Ordnungsstatistiken (P90 bis Maximum von je 90 Aufrufen) und Anteil unter 2\,s über die drei Läufe (270 Aufrufe) mit
    95-\%-Intervall nach Clopper-Pearson}
  \label{tab:kaltstart-ki}
  \footnotesize
  \setlength{\tabcolsep}{4pt}
  \begin{tabularx}{\textwidth}{@{}>{\RaggedRight\arraybackslash}X r r r r r@{}}
    \hline
    Bedingung & P95, Lauf 1 & P95, Lauf 2 & P95, Lauf 3 & unter 2\,s & Intervall \\
    \hline
    Ausgang, verteilt & 3,12 bis 3,31\,s & 3,22 bis 5,03\,s & 2,92 bis 5,18\,s & 3 & 0,2 bis 3,2\,\% \\
    Pool 40, verteilt & 358 bis 555\,ms & 350\,ms bis 1,23\,s & 406\,ms bis 1,37\,s & 270 & 98,6 bis 100\,\% \\
    Pool 0, verteilt & 5,70 bis 7,59\,s & 6,06 bis 23,5\,s & 6,32 bis 9,65\,s & 99 & 30,9 bis 42,7\,\% \\
    Ausgang, Schub & 14,1 bis 20,2\,s & 12,7 bis 21,8\,s & 15,3 bis 24,3\,s & 0 & 0 bis 1,4\,\% \\
    Pool 40, Schub & 6,45 bis 10,8\,s & 6,89 bis 11,0\,s & 6,85 bis 8,72\,s & 119 & 38,1 bis 50,2\,\% \\
    Pool 0, Schub & 9,90 bis 12,6\,s & 10,4 bis 14,5\,s & 10,9 bis 13,9\,s & 0 & 0 bis 1,4\,\% \\
    \hline
  \end{tabularx}
\end{table}

\begin{table}[H]
  \caption{95-\%-Bootstrap-Intervalle des Overheads in Millisekunden je Lauf vom 29.\ Oktober, Blöcke von 1.000
    aufeinanderfolgenden Aufrufen, 1.000 Ziehungen}
  \label{tab:konfidenz}
  \small
  \begin{tabularx}{\textwidth}{@{}>{\RaggedRight\arraybackslash}X r r r r@{}}
    \hline
    Lauf & P50 & P90 & P95 & P99 \\
    \hline
    P1, Prototyp & 143 bis 148 & 227 bis 249 & 284 bis 296 & 354 bis 371 \\
    P2, Prototyp & 138 bis 141 & 177 bis 186 & 208 bis 233 & 301 bis 327 \\
    P3, Prototyp & 141 bis 147 & 194 bis 217 & 241 bis 267 & 326 bis 345 \\
    K1, Ausgangszustand & 305 bis 312 & 474 bis 530 & 676 bis 875 & 3.204 bis 5.045 \\
    K2, Ausgangszustand & 308 bis 316 & 490 bis 541 & 690 bis 857 & 3.061 bis 4.628 \\
    K3, Ausgangszustand & 298 bis 302 & 445 bis 494 & 670 bis 814 & 3.066 bis 4.483 \\
    \hline
  \end{tabularx}
\end{table}


\begin{table}[H]
  \caption{Zuordnung der Messgrößen zu Skripten und Läufen}
  \label{tab:messgroessen}
  \small
  \begin{tabularx}{\textwidth}{@{}l >{\RaggedRight\arraybackslash}X >{\RaggedRight\arraybackslash}X@{}}
    \hline
    Messgröße & Quelle & Ausgangszustand \\
    \hline
    S1 Kaltstart & Kaltstart-Skript, Metrik \texttt{cold\_start\_duration}, Läufe 00 bis 08, Einmalmodus auf hyper-int-otc, Pool 40 und Pool 0
      je dreimal am 29.\ Oktober & Einmalmodus, je dreimal am 29.\ Oktober \\
    S2 warme Latenz & Overhead im Lastprofil, drei Läufe am 29.\ Oktober, Zerlegung über Traces & drei Läufe am 29.\ Oktober \\
    S3 Speicher & Containerstatistiken je Pod unter dem Lastprofil, Ruhezustand des Pools, 29.\ und 30.\ Oktober &
      Containerstatistiken je Pod unter dem Lastprofil \\
    S4 Durchsatz & Stufenlauf mit steigender Ankunftsrate in beiden Namespaces, 30.\ Oktober & ebenso \\
    Fehlerrate & fehlgeschlagene Aufrufe je Lauf, ohne clientseitige Netzfehler & drei Läufe am 29.\ Oktober \\
    Agentenprofil & k6 gegen den Agent-Endpunkt, 2.\ November: warm, Erstaufruf, Größe, ausgehende Anfragen, Parallelität & kein Endpunkt \\
    \hline
  \end{tabularx}
\end{table}

\begin{table}[H]
  \caption{Kaltstartläufe mit dem Kaltstart-Skript, k6, Stand des Prototyps je Lauf}
  \label{tab:kaltstart-verlauf}
  \footnotesize
  \begin{tabularx}{\textwidth}{@{}l >{\RaggedRight\arraybackslash}X r r r r@{}}
    \hline
    Lauf & Stand des Prototyps & Kalt Med. & Kalt P95 & Kalt Max & Warm Med. \\
    \hline
    06.10. & Pool 3, Wartezeit 5\,s, Sidecar. 10 Mandanten, 10 \acsp{VU} & 7,44\,s & 8,84\,s & 9,22\,s & 96\,ms \\
    06.10. & ohne Sidecar, Wartezeit 20\,s, Pool 10. 10 Mandanten, 10 \acsp{VU} & 192\,ms & 215\,ms & 216\,ms & 102\,ms \\
    00 & 10 Mandanten, 40 \acsp{VU}, 5 Aufrufe verloren & 594\,ms & 55,0\,s & 55,0\,s & 142\,ms \\
    01 & 33 Mandanten, 40 \acsp{VU}, Pool 10 & 8,12\,s & 15,5\,s & 21,3\,s & 117\,ms \\
    02 & Pool 40 & 635\,ms & 1,07\,s & 1,22\,s & 130\,ms \\
    03 & Skalierungsschwelle 0,25 & 747\,ms & 1,12\,s & 1,14\,s & 128\,ms \\
    04 & Skalierungsschwelle 0,75 & 613\,ms & 964\,ms & 1,11\,s & 141\,ms \\
    05 & Start der \acsp{VU} über 10\,s verteilt & 350\,ms & 479\,ms & 801\,ms & 124\,ms \\
    06 & ruhige Mandanten, 15\,min, 12 Aufrufe verloren & 189\,ms & 367\,ms & 55,2\,s & -- \\
    07 & nach Neuzuweisung toter Worker & 201\,ms & 382\,ms & 2,41\,s & -- \\
    08 & Slot bei der Bestätigung frei & 396\,ms & 482\,ms & 677\,ms & 178\,ms \\
    \hline
  \end{tabularx}
\end{table}

\begin{table}[H]
  \caption{Zerlegung eines Pod-Starts bis zur ersten Ausführung, 6.\ Oktober, Image auf dem Knoten}
  \label{tab:podstart}
  \small
  \begin{tabularx}{\textwidth}{@{}>{\RaggedRight\arraybackslash}X r@{}}
    \hline
    Schritt & Dauer \\
    \hline
    Pod anlegen und Scheduling & 0,2\,s \\
    Container-Start & 0,6\,s \\
    Prozessstart bis zum ersten Verbindungsversuch & 0,7\,s \\
    Verbindung zu NATS vom Cluster-Netz blockiert & 1,6\,s \\
    Anmeldung, Zuweisung, erste Ausführung & 0,3\,s \\
    Summe, fünf von sieben Pods & 3,3 bis 3,6\,s \\
    Ausreißer: Image-Pull auf einem Knoten ohne Worker & 7,8\,s \\
    \hline
  \end{tabularx}
\end{table}

\begin{table}[H]
  \caption{Ressourcen auf hyper-int-otc während der Kaltstartmessung, Pod-Ebene unter gVisor}
  \label{tab:ressourcen-vergleich}
  \footnotesize
  \begin{tabularx}{\textwidth}{@{}>{\RaggedRight\arraybackslash}X r r r@{}}
    \hline
     & Ausgang & Pool 40 & Pool 0 \\
    \hline
    Neue Runtime-Pods im Fenster & 191 & 188 & 188 \\
    Speicher-Spitze je Pod, Median / P90 / Max, MiB & 89,7 / 92,6 / 104 & 45,8 / 48,8 / 53 & 45,8 / 48,2 / 52 \\
    \acs{CPU} je Pod über die Lebensdauer, Kerne & 0,017 & 0,020 & 0,022 \\
    Lebensdauer eines neuen Pods, Median & 78\,s & 190\,s & 88\,s \\
    Requests \acs{CPU} / Speicher & 10\,m + 25\,m / 128\,Mi & 10\,m / 32\,Mi & 10\,m / 32\,Mi \\
    Limits \acs{CPU} / Speicher & 100\,m / 512\,Mi & -- / 512\,Mi & -- / 512\,Mi \\
    Runtime-Pods, Spitze & 162 & 216 & 185 \\
    Runtime-Speicher, Spitze & 13,7\,GiB & 9,7\,GiB & 8,2\,GiB \\
    Runtime-Speicher vor dem Test & rund 0 & 1,9\,GiB & 0,3\,GiB \\
    Steuerung und Routing, Pods & 7 & 2 bis 3 & 2 bis 3 \\
    Steuerung und Routing, Speicher & 1,35 bis 1,38\,GiB & 113 bis 167\,MiB & 116 bis 167\,MiB \\
    Steuerung und Routing, \acs{CPU} & bis 0,14 Kerne & bis 0,05 Kerne & bis 0,04 Kerne \\
    \hline
  \end{tabularx}
\end{table}

\begin{table}[H]
  \caption{Speicher auf hyper-int-otc nach der Trennung in Transport und Sandbox, 26.\ Oktober}
  \label{tab:ressourcen-trennung}
  \footnotesize
  \begin{tabularx}{\textwidth}{@{}>{\RaggedRight\arraybackslash}X r r@{}}
    \hline
     & Pool 40 & Pool 0 \\
    \hline
    Speicher-Spitze je neuem Pod, Median / P90 / Max, MiB & 58,9 / 64,2 / 68,5 & 59,9 / 63,3 / 67,2 \\
    Control Plane und Gateway, Spitze & 161\,MiB & -- \\
    \hline
  \end{tabularx}
\end{table}

\begin{table}[H]
  \caption{Fehlgeschlagene oder verspätete Aufrufe, ihre Ursache und die Korrektur}
  \label{tab:zustellung}
  \footnotesize
  \begin{tabularx}{\textwidth}{@{}l >{\RaggedRight\arraybackslash}X >{\RaggedRight\arraybackslash}X@{}}
    \hline
    Lauf & Befund und Ursache & Korrektur \\
    \hline
    00 & 5 Aufrufe im Timeout. Ein Pull an einen frisch angelegten Consumer mit drei Repliken
      ging verloren, bevor dieser einen Leader hatte. & Nach dem Anlegen bis zu 2\,s auf den Leader
      warten, einen Pull ohne Heartbeat neu aufbauen. \\
    06 & 12 Aufrufe im Timeout, ausgeführt erst 133 bis 173\,s später mit dem nächsten Worker.
      Der Worker hatte sich 0,01 bis 1,94\,s vor dem Auftrag beendet. Die Control Plane sah den
      Mandanten noch als zugewiesen und bestätigte den Auftrag, ohne dass jemand den Consumer las. & Nach dem Abräumen eines toten Workers Mandanten mit wartenden
      Aufträgen sofort neu zuweisen, auf hyper-int-otc zusätzlich mit Wiederholung nach 1 bis
      16\,s, weil die erste Neuzuweisung einmal scheiterte. Seit dem 26.\ Oktober zieht die Control
      Plane Mandanten ohne Aufträge selbst zurück, bevor ein Worker sich beendet. \\
    19 & 504 ab der ersten Minute auf hyper-int-otc. Das NATS-Konto erlaubte 100 Consumer je
      Stream, ab dem 97.\ aktiven Mandanten bekam keiner einen Worker. & Limit auf 1.000
      angehoben, Lauf aus dem Vergleich ausgeschlossen. \\
    22--29 & 50 bis 109 Zuweisungen je Lauf scheitern, weil Antworten der JetStream-\acs{API} auf
      einzelnen langlebigen Verbindungen ausbleiben: In 276 von 276 Fällen war die Verbindung
      verbunden, ein PING wurde in höchstens 3\,ms beantwortet, und in den zwei Sekunden kam kein
      Byte an, während frische Verbindungen aus einem Prüf-Pod dieselben Consumer ohne einen
      einzigen Fehler lasen. Serverseitige Ursache nicht belegt, an die Plattform gemeldet. &
      Bei genau diesem Muster vor dem nächsten Versuch neu verbinden. Lauf 30: keine
      gescheiterte Zuweisung mehr, 99.\ Perzentil der Zuweisung 2,1 statt 6,0\,s, und 39 der 62
      Reconnects landeten auf demselben Server. \\
    27 & 81 Aufrufe im Timeout nach einem Umbau, der den Consumer direkt zog: Bei fehlendem
      Heartbeat verwarf der Code den Pull und legte den Consumer an, der schon bestand. Der Server
      stellte die wartende Nachricht in den verworfenen Pull zu. & Umbau zurückgenommen. \\
    30 & 1 Aufruf im Timeout. Der Worker hielt einen dauerhaften Pull offen, die Nachricht wurde
      in ihn zugestellt, während der Prozess sich wegen Inaktivität beendete. & Ein Pull je
      Nachricht, Beenden und Drain nur zwischen zwei Pulls, noch ohne Lastlauf. \\
    \hline
  \end{tabularx}
\end{table}

\begin{table}[H]
  \caption{SandboxEval, 51 Eigenschaften für JavaScript angepasst, 21.\ Oktober 2026. Eine Antwort heißt
    nicht, dass der Zugriff auf den echten Host gelang, im Webhook-Profil sind die Werte leer oder nachgebildet}
  \label{tab:sandboxeval}
  \small
  \begin{tabularx}{\textwidth}{@{}>{\RaggedRight\arraybackslash}X r r r@{}}
    \hline
     & Ausgangszustand & Prototyp, Webhook & Prototyp, Agent \\
    \hline
    Fälle, die antworteten (von 51) & 31 & 29 & 2 \\
    davon mit Information über den echten Host & 12 & 0 & 0 \\
    davon mit Verbindung zu externen Zielen & 5 & 0 & 0 \\
    \hline
  \end{tabularx}
\end{table}

\begin{table}[H]
  \caption{Konfigurationsvergleich der Pod-Spezifikation für M3, Container mit Gastcode}
  \label{tab:m3}
  \footnotesize
  \begin{tabularx}{\textwidth}{@{}l >{\RaggedRight\arraybackslash}X >{\RaggedRight\arraybackslash}X >{\RaggedRight\arraybackslash}X@{}}
    \hline
    Feld & Ausgangszustand & Prototyp & Prototyp gegenüber dem Ausgangszustand \\
    \hline
    \texttt{runtimeClassName} & \texttt{gvisor} & \texttt{gvisor} & gleich \\
    \texttt{allowPrivilegeEscalation} & \texttt{false} & \texttt{false} & gleich \\
    \texttt{capabilities.drop} & \texttt{ALL} & \texttt{ALL} & gleich \\
    \texttt{seccompProfile} & \texttt{RuntimeDefault} & \texttt{RuntimeDefault} & gleich \\
    \texttt{automountServiceAccountToken} & \texttt{false} & \texttt{false} & gleich \\
    \texttt{hostNetwork}, \texttt{hostPID}, \texttt{hostIPC} & nicht gesetzt & nicht gesetzt & gleich \\
    \texttt{enableServiceLinks} & \texttt{false} & \texttt{true} & lockerer, Umgebung im Pod, nicht Kernel-Oberfläche \\
    \texttt{readOnlyRootFilesystem} & nicht gesetzt & \texttt{true} & strenger \\
    \texttt{runAsNonRoot} & nicht gesetzt & \texttt{true} & strenger \\
    \hline
  \end{tabularx}
\end{table}

\begin{table}[H]
    \caption{Knotenentzug (M5): Lücke ohne Ausführungsfähigkeit und längste kurze Latenz je Lauf, ab
    Beginn der Löschung, am 28.\ Oktober als längster Abstand zweier Antworten des Probers}
  \label{tab:m5}
  \footnotesize
  \setlength{\tabcolsep}{4pt}
  \begin{tabularx}{\textwidth}{@{}l r r r r r >{\RaggedRight\arraybackslash}X@{}}
    \hline
        Lauf & Getroffen & Lücke getroffen & Lücke übrige & Latenz max. & Fehler & Lange Aufrufe \\
    \hline
    \multicolumn{7}{@{}l}{\emph{Prototyp, mit langen Aufrufen}} \\
    1 & 2 & 1,11\,s & 1,06\,s & 1.269\,ms & 0 & 8 von 8 \\
    2 & 2 & 1,20\,s & 1,13\,s & 1.291\,ms & 0 & 8 von 8 \\
    3 & 3 & 1,37\,s & 0,77\,s & 1.493\,ms & 0 & 8 von 8 \\
    4 & 2 & 0,76\,s & 0,98\,s & 1.261\,ms & 0 & 8 von 8 \\
    \multicolumn{7}{@{}l}{\emph{Prototyp, ohne lange Aufrufe}} \\
    5 & 2 & 1,12\,s & 1,11\,s & 1.483\,ms & 0 & -- \\
    6 & 1 & 0,61\,s & 1,45\,s & 1.303\,ms & 0 & -- \\
    7 & 1 & 0,66\,s & 0,59\,s & 1.243\,ms & 0 & -- \\
    8 & 2 & 1,33\,s & 1,12\,s & 1.271\,ms & 0 & -- \\
        \multicolumn{7}{@{}l}{\emph{Ausgangszustand}} \\
    1 & 1 (3 Pods) & 1,9\,s & -- & 3.000\,ms & 0 & 1 von 1 (26,2\,s) \\
    \multicolumn{7}{@{}l}{\emph{Prototyp, Prober 100\,ms, 29.\ Oktober}} \\
    P1 bis P3 & je 2 & 0,19 bis 0,25\,s & -- & -- & 0 bis 1 (Client) & 24 von 24 (25,2\,s), 48 von 48 Ausführungen genau einmal \\
    \multicolumn{7}{@{}l}{\emph{Ausgangszustand, Prober 100\,ms, 29.\ Oktober}} \\
    K1 bis K3 & je 2 & 2,06 bis 2,68\,s & -- & -- & 0 bis 4 & 24 von 24 (25,3 bis 25,5\,s) \\
    \multicolumn{7}{@{}l}{\emph{Drain mit 40 Mandanten, 29.\ Oktober, Prober 200\,ms}} \\
    Prototyp & 6 & $\le$ 0,56\,s & -- & -- & 2 (Client) & 40 von 40, 80 Ausführungen genau einmal \\
    Knative & 6 & 2,9 bis 5,15\,s & -- & -- & 7 & 39 von 40 \\
    \multicolumn{7}{@{}l}{\emph{Prototyp, Prober 100\,ms, 28.\ Oktober}} \\
    P1 & 3 & 0,22 bis 0,30\,s & -- & -- & 0 & 3 neu ausgeführt (33,4 bis 33,7\,s) \\
    P2 & 4 & 0,24 bis 0,34\,s & -- & -- & 0 & 2 neu ausgeführt (36,5\,s) \\
    P3 & 2 & 0,27 bis 0,42\,s & -- & -- & 0 & 1 neu ausgeführt (35,5\,s) \\
    P4, nach dem Fix & 5 & 0,21 bis 0,27\,s & -- & -- & 0 & 8 von 8 (25,2\,s) \\
    \multicolumn{7}{@{}l}{\emph{Ausgangszustand, Prober 100\,ms, 28.\ Oktober}} \\
    A1 & 1 & 2,42\,s & -- & -- & 0 & 1 von 1 (25,3\,s) \\
    A2 bis A4 & je 3 & 2,10 bis 2,69\,s & -- & -- & 0 & 9 von 9 (25,3\,s) \\

    \hline
  \end{tabularx}
\end{table}



\section{Prüffallkatalog}
\label{sec:prueffallkatalog}

\begin{table}[H]
  \caption{Prüffälle zur Zustandsfreiheit und zur Durchsetzung der Fähigkeitsbeschränkung}
  \label{tab:prueffaelle}
  \small
  \begin{tabularx}{\textwidth}{@{}l >{\RaggedRight\arraybackslash}X l l@{}}
    \hline
    Nr. & Prüffall & Kriterium & Quelle \\
    \hline
    T1 & Ein über \texttt{setTimeout} registrierter Callback überdauert das Ende seiner
         Ausführung und greift auf Daten einer späteren Ausführung zu. & M2 & beobachtet \\
    T2 & Eine in Ausführung $n$ gesetzte globale Variable ist in Ausführung $n+1$ lesbar. &
         M2 & abgeleitet \\
    T3 & Eine Veränderung eines eingebauten Prototyps oder des Module Cache wirkt in
         Ausführung $n+1$ fort. & M2 & abgeleitet \\
    T4 & Eine in Ausführung $n$ abgelegte Datei ist in Ausführung $n+1$ lesbar. & M2 &
         abgeleitet \\
    T5 & Eine in Ausführung $n$ geöffnete Netzwerkverbindung ist in Ausführung $n+1$ nutzbar.
         & M2 & abgeleitet \\
    T6 & Eine Ausführung erschöpft den Speicher und macht die Instanz für folgende
         Ausführungen unbrauchbar. & M2 & beobachtet \\
    T7 & Ausgehender Netzwerkverkehr wird an dem vermittelten Binding vorbei aufgebaut. &
         M4 & abgeleitet \\
        T8 & Auf das Dateisystem wird an dem vermittelten Binding vorbei zugegriffen. & M4 &
         abgeleitet \\
        T9 & Eine Ausführung überschreibt jede ihr sichtbare Kopie ihrer Identität und handelt
         danach im Namen eines anderen Mandanten. & M4 & abgeleitet \\
    T10 & Code mit der Netz- und Prozesssicht der Sandbox verbindet sich ohne gültiges Token mit NATS oder weist dem
         eigenen Transport einen Mandanten zu. & M1 & abgeleitet \\
    \hline
  \end{tabularx}
\end{table}
